\chapter{Introduction}

Backend development is the work of building systems that other systems depend on.
A backend receives requests, validates them, applies business rules, changes state, and returns results.

Most teams do not fail because they cannot write endpoints.
They fail because the backend becomes hard to change.
This happens when interfaces drift, data models become inconsistent, and operational reality is ignored.

\section{A reference system}
Throughout the book, assume a simple product:
a \textbf{Task API} that supports projects, tasks, and activity logs.
It is enough to discuss real constraints:
authentication, pagination, background work, indexing, and deployments.

\section{What ``good'' looks like}
A backend is in good shape when:
\begin{itemize}[nosep]
  \item changes are local and predictable
  \item failures are visible and recoverable
  \item performance is measured, not guessed
  \item security does not rely on developer memory
  \item operations are boring
\end{itemize}

\section{Scope and non-goals}
This book focuses on production systems: APIs, data, and operations.
It does not cover UI concerns or low-level kernel tuning.
If you keep the system simple and observable, most hard problems become manageable.

\section{A minimal API surface}
A reasonable starting API for the Task system:
\begin{itemize}[nosep]
  \item \code{GET /projects}
  \item \code{POST /projects}
  \item \code{GET /projects/:id}
  \item \code{POST /projects/:id/tasks}
  \item \code{PATCH /tasks/:id}
  \item \code{GET /tasks?project=:id}
\end{itemize}

\section{Quality attributes}
Backends are judged by qualities, not just features.
These attributes are the baseline for production systems.

\begin{tabularx}{\textwidth}{@{}lX@{}}
\toprule
Attribute & What it means in practice \\
\midrule
Reliability & Requests succeed or fail predictably, with safe retries \\
Performance & Latency stays within a known budget at peak load \\
Security & Access is controlled, secrets are protected, audit trails exist \\
Scalability & Capacity grows without redesigning the entire system \\
Maintainability & Changes are small, testable, and reversible \\
\bottomrule
\end{tabularx}

\section{The operating model}
Every backend is an operations system.
If you cannot deploy, observe, and recover safely, features do not matter.
Make operational tasks part of the design:
\begin{itemize}[nosep]
  \item define health checks and readiness signals
  \item track error budgets and latency budgets
  \item build runbooks for common failures
\end{itemize}

\section{A small data model}
The Task system is intentionally small:
\begin{itemize}[nosep]
  \item \textbf{Project}: id, name, owner
  \item \textbf{Task}: id, project id, title, status, due date
  \item \textbf{Activity}: id, task id, action, timestamp
\end{itemize}

This is enough to discuss ownership, filtering, pagination, and background work.

\section{Change as the primary constraint}
The hardest part of backend work is change.
You will add fields, remove endpoints, move data, and split services.
A good backend is one where these changes are possible without large rewrites.

\begin{patternbox}{Pattern: treat interfaces as products}
The API is not a technical detail.
It is a contract with clients, future developers, and your operations team.
\end{patternbox}

\section{Common failure modes}
Teams often fall into the same traps:
\begin{itemize}[nosep]
  \item business logic hidden in controllers and frameworks
  \item inconsistent error formats across endpoints
  \item databases without constraints or indexes
  \item no measurable latency or error budgets
  \item ad-hoc security rules
\end{itemize}

\section{Working style}
This book assumes you work in small, focused changes:
\begin{itemize}[nosep]
  \item keep endpoints small and composable
  \item move complex logic into services
  \item test the critical path first
  \item monitor before you optimize
\end{itemize}

\section{Who this is for}
This is for engineers who ship and operate backends:
\begin{itemize}[nosep]
  \item product teams running APIs in production
  \item platform teams enabling internal services
  \item engineers migrating from prototypes to durable systems
\end{itemize}

\section{System boundaries}
Backends are not just servers.
They include storage, queues, caches, and the tooling around them.
Draw a boundary around what you own and what you depend on.
Every dependency outside that boundary needs monitoring and a fallback plan.

\section{A sample request}
A common request in the Task API is creating a task.
The response is intentionally simple:

\begin{lstlisting}[language=JSON]
{
  "id": "task_123",
  "project_id": "proj_7",
  "title": "Write docs",
  "status": "open",
  "created_at": "2026-02-18T10:00:00Z"
}
\end{lstlisting}

\section{Operational baseline}
Every backend should define a baseline operational posture:
\begin{itemize}[nosep]
  \item health checks and readiness endpoints
  \item log correlation identifiers
  \item a dashboard for latency, errors, and traffic
  \item alerts tied to user-visible symptoms
\end{itemize}

\section{Documentation artifacts}
Keep lightweight documentation close to the code:
\begin{itemize}[nosep]
  \item an API reference or OpenAPI spec
  \item runbooks for the top three incidents
  \item a short architecture overview for new engineers
\end{itemize}

\section{What a backend is responsible for}
Backends are responsible for more than serving HTTP:
\begin{itemize}[nosep]
  \item enforcing business rules consistently
  \item safeguarding data integrity
  \item managing concurrency and coordination
  \item providing stable contracts to clients
  \item operating safely under load
\end{itemize}

\section{How backends fail in practice}
Backends usually fail slowly before they fail completely.
Latency creeps up, error rates climb, and on-call becomes reactive.
If you monitor latency, errors, and saturation, you can intervene early.

\section{A short architecture sketch}
The Task API can be implemented with a simple architecture:
\begin{itemize}[nosep]
  \item an API service for HTTP requests
  \item a relational database for tasks and projects
  \item a queue for activity events and background jobs
  \item a cache for hot reads
\end{itemize}

\section{Development lifecycle}
Production systems require a lifecycle:
\begin{itemize}[nosep]
  \item local development with realistic data
  \item staged environments for integration tests
  \item production deploys with rollback paths
  \item post-deploy monitoring and validation
\end{itemize}

\section{Tradeoffs are expected}
There is no perfect design.
You choose tradeoffs based on expected load, team size, and product risk.
Make those tradeoffs explicit and document them.

\section{From prototype to product}
Most backends start as a prototype.
The difference between prototype and product is the ability to change safely.
Add constraints, monitoring, and tests before you add more endpoints.

\section{Operational constraints}
Production systems operate under constraints:
\begin{itemize}[nosep]
  \item limited error budgets
  \item fixed database capacity
  \item on-call fatigue
  \item unpredictable traffic patterns
\end{itemize}

\section{Interfaces beyond HTTP}
Your backend interfaces include more than HTTP endpoints.
Database schemas, queues, and background job payloads are also contracts.
Treat those contracts with the same care as API endpoints.

\section{Release quality gates}
Define a minimal set of gates before each deploy:
\begin{itemize}[nosep]
  \item unit and integration tests
  \item database migration checks
  \item smoke tests on staging
  \item rollback validation
\end{itemize}

\section{Who depends on you}
Clients are not just user interfaces.
They include other services, data pipelines, and internal tooling.
The more dependents you have, the more you need stable interfaces.

\section{Reference workflow}
A minimal workflow for backend changes:
\begin{itemize}[nosep]
  \item write or update the API contract
  \item implement the change in the service layer
  \item add tests for the critical path
  \item update migrations and run locally
  \item deploy to staging and validate metrics
\end{itemize}

\section{Baseline metrics}
Start with a small, stable set of metrics:
\begin{itemize}[nosep]
  \item request rate per endpoint
  \item error rate per endpoint
  \item latency p50, p95, p99
  \item database query latency
  \item queue depth for background work
\end{itemize}

\section{Developer experience}
If the development loop is slow, engineers will skip tests and avoid refactors.
Invest in a fast local setup, representative test data, and deterministic tooling.

\section{Case study: growth without chaos}
The Task API starts with a single service and a single database.
As usage grows, new requirements appear: search, analytics, and integrations.
If the API contract and data model are stable, these changes are incremental.
If they are ad-hoc, every new feature becomes a migration project.

The lesson is simple: design for change early.
The first version does not need to be perfect, but it must be consistent.

\section{Configuration and environments}
Backends live in multiple environments: local, staging, and production.
Configuration must be external to code and validated on startup.
If a service cannot start without configuration, fail fast and clearly.

\section{Feature flags}
Feature flags are a safe way to roll out risky changes.
They allow you to test in production with a small subset of traffic.
Use flags sparingly and remove them after the rollout.

\section{Data contracts}
Data is a contract just like the API.
When you add or change a field, you are changing the contract.
Document those changes and plan migrations explicitly.

\section{Incident response}
Expect incidents and rehearse your response.
A good incident response includes:
\begin{itemize}[nosep]
  \item clear ownership and on-call rotation
  \item runbooks for top failure modes
  \item post-incident reviews with action items
\end{itemize}

\section{Scaling the team}
As teams grow, communication becomes the bottleneck.
Stable interfaces, clear ownership, and shared standards keep teams aligned.

\section{Operational posture}
A stable backend has a known operational posture.
That posture includes how you deploy, how you observe, and how you recover.
If these are undocumented, the system will only be safe in the heads of a few engineers.

\section{Production constraints you cannot ignore}
Every production system has constraints that shape design:
\begin{itemize}[nosep]
  \item limited database connections
  \item external API rate limits
  \item on-call hours and handoff processes
  \item compliance and audit requirements
\end{itemize}

\section{API evolution mindset}
Assume the API will evolve continuously.
Design so that adding fields and new endpoints is routine.
Treat breaking changes as exceptional events with a clear process.

\section{Example ownership map}
Ownership keeps change safe.
A simple ownership map for the Task system might be:
\begin{itemize}[nosep]
  \item Projects: core backend team
  \item Tasks: core backend team
  \item Activity logs: platform team
  \item Search and analytics: data team
\end{itemize}

\section{Release reliability}
Reliability is not just runtime behavior.
It includes the release process: how you migrate data, roll out changes, and rollback.
If releases are unsafe, production will be unsafe too.

\section{Assumptions and load profile}
Every backend makes assumptions about scale and shape of traffic.
Write them down, even if they are rough:
\begin{itemize}[nosep]
  \item expected requests per second at peak
  \item typical request sizes and response sizes
  \item daily batch windows or traffic spikes
  \item read-to-write ratio for core endpoints
\end{itemize}

These assumptions guide indexing, caching, and capacity planning.
When they change, the system needs to change with them.

\section{Security baseline}
Security is part of the contract.
A minimal baseline includes:
\begin{itemize}[nosep]
  \item explicit authentication and authorization checks
  \item secrets managed outside the codebase
  \item encrypted transport and sensitive data handling
  \item an audit trail for sensitive operations
\end{itemize}

If these are implicit, they will be inconsistent.

\section{Cost awareness}
Cost is a constraint like latency and availability.
Small decisions compound into large bills.
Start with a simple discipline:
\begin{itemize}[nosep]
  \item know the cost drivers (database, queue, storage, egress)
  \item track usage at the service level
  \item set budgets for non-critical workloads
\end{itemize}

\section{Dependency management}
Every backend depends on other systems.
Make those dependencies explicit:
\begin{itemize}[nosep]
  \item document upstream and downstream services
  \item define timeouts and retry policies per dependency
  \item record the expected availability and failure modes
\end{itemize}

\section{Architecture decisions}
Some choices are durable: database, messaging system, schema strategy.
Capture them in short Architecture Decision Records.
The goal is not documentation volume, but traceability:
future engineers should see why a decision was made.

\section{Readiness review}
Before calling a system production-ready, ask:
\begin{itemize}[nosep]
  \item Can we deploy without downtime?
  \item Can we restore data from backup?
  \item Do we know our p95 latency for key endpoints?
  \item Can we detect a silent failure within minutes?
\end{itemize}

\section{Operational drills}
Backups are only real after a restore.
Run at least one restore drill per quarter and record the outcome.
Treat drills as a feature: a backend that cannot be restored is incomplete.

\section{Service catalog}
As your system grows, keep a service catalog.
It should list owners, on-call rotations, dependencies, and critical endpoints.
This prevents knowledge from living only in people's heads.

\section{Risk register}
Every backend has known risks.
Capture the top three risks and update them as the system changes:
\begin{itemize}[nosep]
  \item single database instance
  \item weak observability for background jobs
  \item ad-hoc access to production data
\end{itemize}

\section{Definition of done}
Feature work is not done until it is safe to operate.
A minimal definition of done includes:
\begin{itemize}[nosep]
  \item tests for the critical path
  \item metrics and dashboards updated
  \item runbook notes for new failure modes
\end{itemize}

\section{When to split services}
Splitting too early adds complexity.
Split when a domain is stable, owned by another team, and scaling needs are clear.
If those conditions are not met, a well-structured modular monolith is safer.

\section{Operating principles}
Good backends follow a small set of principles:
\begin{itemize}[nosep]
  \item keep the critical path short
  \item prefer explicit contracts over implicit behavior
  \item make failures visible before they become outages
  \item optimize for reversible change
\end{itemize}

If a design violates these principles, it will be expensive to operate later.

\section{Stakeholder alignment}
Backend work impacts product, data, and operations.
Align early on:
\begin{itemize}[nosep]
  \item data ownership and access policies
  \item SLA expectations from clients and partners
  \item maintenance windows and on-call coverage
\end{itemize}

Alignment is not paperwork; it prevents future incidents.

\section{Example release plan}
A simple release plan for a risky change:
\begin{itemize}[nosep]
  \item ship a backward-compatible schema change
  \item deploy code that reads old and new data
  \item backfill existing records
  \item enable the new behavior behind a feature flag
  \item remove the old field after a full rollout
\end{itemize}

\begin{checklistbox}{Checklist: starting point}
\begin{itemize}[nosep]
  \item Define the core resources and their identifiers.
  \item Decide how errors will be represented.
  \item Write the simplest endpoints that cover the use cases.
  \item Create a baseline dashboard for latency and errors.
\end{itemize}
\end{checklistbox}
